\documentclass[a4paper,11pt]{ltjsarticle}
\usepackage[haranoaji]{luatexja-preset}
\usepackage{amsmath,amssymb}
\usepackage[top=12mm,bottom=10mm,left=14mm,right=14mm]{geometry}
\usepackage{array}
\usepackage{tikz}
\usetikzlibrary{calc}
\pagestyle{empty}
\setlength{\parindent}{0pt}
\newlength{\qcol}\setlength{\qcol}{0.47\textwidth}
\newlength{\qgap}
\newlength{\anscolw}\setlength{\anscolw}{0.30\textwidth}
\newlength{\ansh}
\newcommand{\Q}[2]{\parbox[t]{\qcol}{\raggedright (#1)\hspace{0.6em}$\displaystyle #2$}}
\newcommand{\QR}[4]{\Q{#1}{#2}\hfill\Q{#3}{#4}\par\vspace{\qgap}}
\newcommand{\anscell}[1]{\rule[-\ansdrop]{0pt}{\ansh}(#1)}
\newlength{\ansdrop}

\begin{document}
{\large\bfseries 計算問題　3}\hfill 名前(\underline{\hspace{32mm}})\par\vspace{3mm}
■（1）〜（16）の計算をしなさい。（17）〜（20）は方程式を解きなさい。\par\vspace{3mm}
\setlength{\qgap}{5.4mm}
\QR{1}{(2x+6)+(5x-3)}{2}{(3a+7)-(-2a-6)}
\QR{3}{25b\div(-5)}{4}{(-6)\times3a}
\QR{5}{2+2^2+2}{6}{(-18)\div(-3)}
\QR{7}{-3+(-7)}{8}{(-4)\times(-5)}
\QR{9}{\dfrac{x-2}{3}\times6}{10}{(25x-15)\div5}
\QR{11}{(+3)+(-9)-(-6)}{12}{8\div5\times(-20)}
\QR{13}{(-4)-(-5)}{14}{\dfrac{1}{2}(6x-4)+\dfrac{1}{3}(3x+9)}
\QR{15}{3-(-8)\div\left(-\dfrac{2}{5}\right)}{16}{\dfrac{1}{4}(8a+12)-\dfrac{1}{2}(10a-6)}
\QR{17}{4x+11=8x+7}{18}{x-6=-x}
\QR{19}{0.3x-0.4=2}{20}{\dfrac{5}{6}x+4=\dfrac{1}{2}x}
\vspace{1mm}
\Q{21}{(-a)^{2018}}\par\vspace{3mm}
\setlength{\ansh}{9.5mm}\setlength{\ansdrop}{6mm}%
\begin{center}\begin{tabular}{|p{\anscolw}|p{\anscolw}|p{\anscolw}|}\hline
\anscell{1} & \anscell{2} & \anscell{3} \\\hline
\anscell{4} & \anscell{5} & \anscell{6} \\\hline
\anscell{7} & \anscell{8} & \anscell{9} \\\hline
\anscell{10} & \anscell{11} & \anscell{12} \\\hline
\anscell{13} & \anscell{14} & \anscell{15} \\\hline
\anscell{16} & \anscell{17} & \anscell{18} \\\hline
\anscell{19} & \anscell{20} & \anscell{21} \\\hline
\end{tabular}\end{center}

\end{document}
